% Generated draft: requires mathematical and source review.
\documentclass[11pt]{article}
\usepackage[margin=25mm]{geometry}
\usepackage{fontspec}
\setmainfont{Latin Modern Roman}
\usepackage{amsmath,amssymb,mathtools,unicode-math}
\setmathfont{Latin Modern Math}
\usepackage{enumitem,needspace,xurl,hyperref}
\setlength{\parindent}{0pt}
\setlength{\parskip}{5pt}
\setlength{\emergencystretch}{3em}
\newcommand{\N}{\mathbb{N}}
\newcommand{\Z}{\mathbb{Z}}
\newcommand{\Q}{\mathbb{Q}}
\newcommand{\R}{\mathbb{R}}
\newcommand{\C}{\mathbb{C}}
\newcommand{\F}{\mathbb{F}}
\newcommand{\A}{\mathbb{A}}
\newcommand{\PP}{\mathbb{P}}
\newcommand{\E}{\mathbb{E}}
\newcommand{\eps}{\varepsilon}
\newcommand{\dd}{\,\mathrm{d}}
\begin{document}
\section*{UnsolvedMath ID: 800915}
\textbf{Source problem code:} catalogue entry
\subsection{Definitions and mathematical statement}
Under the customary complex-projective interpretation, an Enriques surface is a smooth connected projective surface \(X/\C\) such that \(H^1(X,\mathcal O_X)=0\), \(\omega_X\not\cong\mathcal O_X\), and \(\omega_X^{\otimes 2}\cong\mathcal O_X\). The supplied excerpt does not explicitly specify the ground field; this interpretation therefore requires confirmation.

Let \(F_g^X\) denote the genus-\(g\) descendant Gromov--Witten potential of \(X\), regarded as a formal generating series in descendant insertion variables and curve-class variables. The total descendant partition function is
\[
Z_X=\exp\!\left(\sum_{g\geq 0}\hbar^{g-1}F_g^X\right).
\]
Here \(\hbar\) is a formal genus parameter. Let \(L_m^X\), for integers \(m\geq -1\), denote the Gromov--Witten Virasoro operators intended by the cited source.

The requested assertion, in this conventional interpretation, is that for every complex projective Enriques surface \(X\),
\[
L_m^X Z_X=0\qquad\text{for every integer }m\geq -1,
\]
as identities of formal generating series, encompassing all genera and curve classes included in the partition function.

This is a statement schema requiring specification: the supplied excerpt gives neither the formulas for \(L_m^X\) nor the normalization of \(Z_X\), descendant-variable conventions, or unstable and degree-zero conventions. Consequently, it does not determine a fully explicit system of equations. The complex ground field and the full descendant interpretation above are stated interpretations, not hypotheses explicitly supplied in the excerpt.
\subsection{Short English statement}
Establish the Gromov--Witten Virasoro constraints for Enriques surfaces. In the customary descendant formulation, every Virasoro operator indexed by \(m\geq -1\) should annihilate the total descendant partition function. The intended ground field and precise operator and generating-function conventions require confirmation from the source.
\subsection{Sources}
\begin{itemize}
\item Local UnsolvedMath registry: permanent ID \(800915\), exact source problem code \texttt{catalogue entry}, title ``AGNES Open Problems in Algebraic Geometry --- Q3.'' The supplied record attributes the problem to Pandharipande, with proposed year 2010, and labels its status as open; this label is supplied metadata, not an independently verified current assessment.
\item Pandharipande, \emph{Open Problems in Enumerative Algebraic Geometry (AGNES)} (2010), source item Q3 (\#2), \url{https://arxiv.org/abs/1004.3259}. The supplied provenance records access on 2026-10-16 and extraction by \texttt{arxiv\_latex}; these are supplied provenance statements, not a claim of independent source access.
\item The supplied excerpt cites the unresolved reference label \texttt{[mp2]} for a study of Enriques-surface Gromov--Witten theory closely related to modular forms. No bibliographic entry or link for this reference is supplied. The excerpt also describes Enriques surfaces as perhaps the most basic variety requiring new techniques to establish the Virasoro constraints; this is the source's qualified assessment in its historical context.
\end{itemize}
\end{document}
